% =====================================================
% 07_discussion.tex — limitations + future work
% =====================================================

\paragraph{C4 is a regime indicator, not a within-regime ordering.}
Our mixed-Hessian predictor $P(a,b)$ has Spearman correlation
$+0.48 / +0.81 / +0.50$ with measured compositional non-linearity
across three architectures. However, dropping the highest-curvature
attribute from each domain collapses bedroom's Spearman to $-0.17$
($n=21$, $p=0.47$) and FFHQ's to $+0.60$ ($n=6$, $p=0.21$). The
single-scalar predictor therefore discriminates \emph{between}
structural and textural curvature regimes but does not order pairs
within the texture regime. A finer per-pixel predictor that
preserves the directional information of $d^2 G(b_a, b_b)$ rather
than collapsing to its norm is the natural next refinement.

\paragraph{Linear boundaries only.}
The Lie-bracket analysis in \S\ref{sec:theory} treats attribute
directions as constant vector fields on $\W$. For learned non-linear
edits --- e.g.\ EditGAN's part-conditioned edit
vectors~\cite{Ling2021EditGAN} or DragGAN's
point-based deformation fields~\cite{Pan2023DragGAN} --- the full Lie
bracket of generally varying vector fields is non-zero and may
dominate the Hessian term. Extending the curvature prediction to such
cases is open.

\paragraph{Diffusion: image-space derivative, not score derivative.}
For Stable Diffusion we measure $\partial^{1,2}\, \mathrm{decode}(x_0)
/\partial(\text{h-space direction})$ via JVP through the DDIM
sampling chain, parallel to the StyleGAN measurements. This is
strictly different from the score-network Hessian
$\nabla^2_x \log p_t(x)$ studied
in~\cite{Stanczuk2024IntrinsicDim, Park2023Riemannian}, which lives
in noisy-latent space and concerns the data distribution's local
geometry. The two views are complementary: ours is generator-side
and operational (predicts edit behaviour), theirs is
distribution-side and structural (counts intrinsic dimension). A
unified treatment is an attractive next step.

\paragraph{Memory ceiling forces FD-of-JVPs on SD.}
Composed JVP through the full 50-step DDIM chain at 512\textsuperscript{2}
OOMs on 8 GB. We therefore compute the second derivative as a
finite-difference of two first-order JVPs at $\pm\epsilon$, which
keeps peak memory at the single-JVP budget. This is the same
construction used by~\cite{Park2023Riemannian} for the diffusion
Jacobian. A larger GPU enables direct composed-JVP and would
reduce the FD truncation error; we report sensitivity-to-$\epsilon$
in Appendix~C.

\paragraph{Forward mode is the right tool here, not a panacea.}
For pushforward and HVP we exploit forward-mode maximally: input
dimension $1$, output dimension $\sim 10^5$ pixels makes forward
$10^3$--$10^4 \times$ faster than reverse-mode at matched accuracy.
For high-dimensional latent-space optimisation (e.g.\ PTI
inversion) reverse mode remains the right choice. We consider
forward mode complementary, not a replacement.

\paragraph{Theoretical assumptions.}
We assume $G$ is $C^2$ at the sampled latents. ReLU non-smoothness
forms a measure-zero set in $\W$; we did not encounter numerical
issues in practice. A fully rigorous statement requires piece-wise
smooth differential geometry --- see supplementary remark.

\paragraph{Beyond GANs and diffusion.}
The differential-geometric view applies, with adaptation, to any
generative model defined by a smooth map from a latent space. Direct
application to VAEs, normalising flows, and consistency-based
generators after suitable smoothing is a clear next step.

\paragraph{Editing claim is descriptive without the head-to-head.}
The cleanest prescriptive use of our metric --- pre-selecting
editing directions by curvature regime --- requires the editing
head-to-head we present in \S\ref{sec:exp-baselines-results}.
In its current form the metric is established as descriptive
(curvature correlates with disentanglement, composition, layer
localisation); the prescriptive payoff is the empirical question
that decides whether the framework changes practice.
